% !PRINTED 164
\weberpage{164}
%<<CONT
identische Verschwinden nicht stattfindet, sondern eine Endgleichung von der Ordnung $mnp$ auftritt, so muss dies auch der
Grad der Endgleichung im allgemeinen Falle sein.

Hiermit ist das \gesperrt{Bezout'sche} Theorem für den Fall von
drei homogenen quaternären Gleichungen bewiesen.

Dass wir uns hier auf drei Gleichungen beschränkt haben,
ist lediglich im Interesse der Einfachheit geschehen. Die Erweiterung, die für den Fall von $n$ Gleichungen mit $n + 1$ homogenen Unbekannten erforderlich ist, ergiebt sich von selbst, und
kann dem Leser überlassen bleiben.

Für unsere allgemeine Aufgabe wollen wir nur noch bemerken, dass sich aus diesen Betrachtungen ergiebt, dass das
Problem der Algebra, sofern es sich auf die Lösung irgend
welcher algebraischer Gleichungen bezieht, damit zurückgeführt
% [corrected 2026-08-07, 8x zoom strip 3/8] The run begins at "ist" -- the query asked where the
% span starts and this settles it: "welcher algebraischer Gleichungen bezieht, damit
% zurückgeführt" is set normally and the very next line is spaced throughout.
\gesperrt{ist auf die Behandlung einer Kette von Gleichungen mit je
einer Unbekannten.}

\begin{center}
\S.\,51.

\gesperrt{Zerlegbare und unzerlegbare Functionen.}
\end{center}

Unter den ganzen rationalen Functionen mehrerer Veränderlichen müssen \gesperrt{zerlegbare und unzerlegbare} unterschieden
werden.

Eine ganze rationale Function von irgend welchen Veränderlichen heisst \gesperrt{zerlegbar}, wenn sie als Product von wenigstens
zwei ganzen rationalen Functionen derselben Veränderlichen darstellbar ist; ist sie nicht als ein solches Product darstellbar, so
heisst sie \gesperrt{unzerlegbar}.

Wenn eine Function zerlegbar ist, so ist der Grad jedes
der Factoren niedriger als der Grad der Function selbst. Eine
\gesperrt{lineare} Function ist also immer unzerlegbar, und \gesperrt{jede ganze
rationale Function lässt sich in eine endliche Anzahl
unzerlegbarer Functionen zerlegen.}

Wir nennen eine ganze Function $W$ durch eine andere $w$
theilbar, wenn eine dritte ganze Function $w'$ existirt, so dass

\[
W = w\,w'
\]

ist. Daraus folgt dann, dass, wenn $U$ eine durch $w$ theilbare
und $V$ eine beliebige ganze Function ist, auch das Product $UV$
%CONT>>
